% =====================================================================
%  Layout, tipografia e titulos. Este arquivo e seu: o latexkit nunca o
%  sobrescreve.
%
%  Ele e carregado DEPOIS de config/metadata.tex, entao e aqui que voce
%  sobrescreve um metadado que precise conter LaTeX:
%
%    \renewcommand{\lkTitle}{O papel do $\alpha$}
% =====================================================================

% --- Mancha de texto -------------------------------------------------
% A memoir calcula margens equilibradas para frente-e-verso: a margem
% interna (lombada) fica menor que a externa.
\setlrmarginsandblock{2.2cm}{2.8cm}{*}
\setulmarginsandblock{2.2cm}{2.6cm}{*}
\checkandfixthelayout

% --- Estilo dos titulos de capitulo ----------------------------------
% A memoir traz varios prontos: default, section, hangnum, companion,
% article, bianchi, bringhurst, brotherton, chappell, crosshead, culver,
% dash, demo2, dowding, ell, ger, komalike, lyhne, madsen, ntglike,
% pedersen, southall, tandh, thatcher, veelo, verville, wilsondob.
\chapterstyle{bianchi}

% --- Cabecalhos e rodapes --------------------------------------------
% Titulo do capitulo na pagina par, titulo da secao na impar, numero da
% pagina na borda externa.
\makepagestyle{lglivro}
\makeevenhead{lglivro}{\thepage}{}{\small\itshape\leftmark}
\makeoddhead{lglivro}{\small\itshape\rightmark}{}{\thepage}
\makeheadrule{lglivro}{\textwidth}{0.4pt}
\pagestyle{lglivro}

% --- Espacamento -----------------------------------------------------
\setlength{\parindent}{1.2em}
\nonzeroparskip
\setlength{\parskip}{0.2em}

% --- Epigrafes -------------------------------------------------------
\setlength{\epigraphwidth}{0.55\textwidth}
\renewcommand{\epigraphflush}{flushright}
\renewcommand{\textflush}{flushright}
\renewcommand{\sourceflush}{flushright}
